% GENERATED FILE. Edit canonical lessons, then run npm run generate:books.
% canonical-source-hash: fnv1a64:a02a21ff2a7d5a64
% canonical-lessons: HI-C416-khema, HI-C416-sulana, HI-C416-basera, HI-C416-guzaarna, HI-C416-bitana

\chapter{Camping Out}
\label{ch:hi-camping-out}
\hlchaptermodality{416}

\begin{chapteropening}
\textbf{By the end of this chapter:} \emph{I can talk about a tent and a roost for the night, settle a child to sleep and say how we spent our days.}

Complete the last lesson of chapter 416: I can talk about a tent and a roost for the night, settle a child to sleep and say how we spent our days.
\end{chapteropening}

\section[\texorpdfstring{\hi{ख़ेमा} (khemā)}{ख़ेमा (khemā)}]{\hi{ख़ेमा} (khemā) — a tent (for camping)}

\label{lesson:HI-C416-khema}



\begin{quote}
Before the new one: say the Hindi for a place to stay, then the Hindi for to send for.
\end{quote}

\subsection*{You'll want to know: \hi{ख़ेमा}}
\textbf{\hi{ख़ेमा}} — \emph{khemā} — \textquotedblleft{}a tent (for camping)\textquotedblright{}.

\begin{quote}
\hi{नदी} \hi{के} \hi{किनारे} \hi{हमारा} \hi{ख़ेमा} \hi{था।}
\end{quote}

\begin{quote}
\emph{nadī ke kināre hamārā khemā thā.}
\end{quote}

\begin{quote}
Our tent was on the river bank.
\end{quote}

Masculine.

\begin{cousinweb}
\textbf{\hi{ख़ेमा}} (\emph{khemā}) is usually traced to Arabic \emph{khayma}, \textquotedblleft{}a tent\textquotedblright{}, which came through Persian; the diphthong has become a plain e.
\end{cousinweb}

\begin{grammarlens}[title={What you've built}]
One more word for getting around.
\end{grammarlens}

\subsection*{Your turn}
\begin{itemize}
\raggedright
  \item \emph{Say it:} \emph{khemā}
  \item \emph{Say it:} \emph{khemā}, once more
  \item \emph{Say it:} \emph{nadī ke kināre hamārā khemā thā.} — Our tent was on the river bank.
  \item \emph{Recall:} say the Hindi for a place to stay, then the Hindi for to send for, then say \emph{khemā} again
\end{itemize}

\begin{tcolorbox}[breakable,colback=teal!4,colframe=teal!35!black,title={Before you move on}]
What does \hi{ख़ेमा} mean? (\textquotedblleft{}A tent (for camping)\textquotedblright{}.) Say it once more.
\end{tcolorbox}
\section[\texorpdfstring{\hi{सुलाना} (sulānā)}{सुलाना (sulānā)}]{\hi{सुलाना} (sulānā) — to put to sleep; to lull (someone) to sleep}

\label{lesson:HI-C416-sulana}



\begin{quote}
Before the new one: say the Hindi for to send for, then the Hindi for a tent (for camping).
\end{quote}

\subsection*{You'll want to know: \hi{सुलाना}}
\textbf{\hi{सुलाना}} — \emph{sulānā} — \textquotedblleft{}to put to sleep; to lull (someone) to sleep\textquotedblright{}.

\begin{quote}
\hi{माँ} \hi{ने} \hi{बच्चे} \hi{को} \hi{ख़ेमे} \hi{में} \hi{सुलाया।}
\end{quote}

\begin{quote}
\emph{māṁ ne bacce ko kheme meṁ sulāyā.}
\end{quote}

\begin{quote}
The mother put the child to sleep in the tent.
\end{quote}

Past, with \hi{ने} (he or she …): \hi{उसने} \hi{बच्ची} \hi{को} \hi{सुलाया}, \emph{usne baccī ko sulāyā}.

\begin{cousinweb}
\textbf{\hi{सुलाना}} (\emph{sulānā}) is a tadbhava of Sanskrit \emph{svāpayati}, \textquotedblleft{}lulls to sleep\textquotedblright{} (CDIAL 13928), which lists Hindi \emph{sulānā}: the causative partner of \hi{सोना}, to sleep.
\end{cousinweb}

\begin{grammarlens}[title={What you've built}]
One more word for getting around.
\end{grammarlens}

\subsection*{Your turn}
\begin{itemize}
\raggedright
  \item \emph{Say it:} \emph{sulānā}
  \item \emph{Say it:} \emph{sulānā}, once more
  \item \emph{Say it:} \emph{māṁ ne bacce ko kheme meṁ sulāyā.} — The mother put the child to sleep in the tent.
  \item \emph{Recall:} say the Hindi for to send for, then the Hindi for a tent (for camping), then say \emph{sulānā} again
\end{itemize}

\begin{tcolorbox}[breakable,colback=teal!4,colframe=teal!35!black,title={Before you move on}]
What does \hi{सुलाना} mean? (\textquotedblleft{}To put to sleep; to lull (someone) to sleep\textquotedblright{}.) Say it once more.
\end{tcolorbox}
\section[\texorpdfstring{\hi{बसेरा} (baserā)}{बसेरा (baserā)}]{\hi{बसेरा} (baserā) — a roost; a night's shelter}

\label{lesson:HI-C416-basera}



\begin{quote}
Before the new one: say the Hindi for a tent (for camping), then the Hindi for to put to sleep.
\end{quote}

\subsection*{You'll want to know: \hi{बसेरा}}
\textbf{\hi{बसेरा}} — \emph{baserā} — \textquotedblleft{}a roost; a night's shelter\textquotedblright{}.

\begin{quote}
\hi{शाम} \hi{को} \hi{चिड़ियाँ} \hi{अपने} \hi{बसेरे} \hi{में} \hi{लौटती} \hi{हैं।}
\end{quote}

\begin{quote}
\emph{śām ko ciṛiyā̃ apne basere meṁ lauṭtī haiṁ.}
\end{quote}

\begin{quote}
In the evening the birds go back to their roost.
\end{quote}

Masculine.

\begin{cousinweb}
\textbf{\hi{बसेरा}} (\emph{baserā}) goes back to \emph{*vāsakara}, \textquotedblleft{}making a stay\textquotedblright{}, a form CDIAL reconstructs (CDIAL 11594), which lists Hindi \emph{baserā}, \textquotedblleft{}a place of staying\textquotedblright{}: a bird's roost, or a traveller's bed for one night.
\end{cousinweb}

\begin{grammarlens}[title={What you've built}]
One more word for getting around.
\end{grammarlens}

\subsection*{Your turn}
\begin{itemize}
\raggedright
  \item \emph{Say it:} \emph{baserā}
  \item \emph{Say it:} \emph{baserā}, once more
  \item \emph{Say it:} \emph{śām ko ciṛiyā̃ apne basere meṁ lauṭtī haiṁ.} — In the evening the birds go back to their roost.
  \item \emph{Recall:} say the Hindi for a tent (for camping), then the Hindi for to put to sleep, then say \emph{baserā} again
\end{itemize}

\begin{tcolorbox}[breakable,colback=teal!4,colframe=teal!35!black,title={Before you move on}]
What does \hi{बसेरा} mean? (\textquotedblleft{}A roost; a night's shelter\textquotedblright{}.) Say it once more.
\end{tcolorbox}
\section[\texorpdfstring{\hi{गुज़ारना} (guzārnā)}{गुज़ारना (guzārnā)}]{\hi{गुज़ारना} (guzārnā) — to spend (time); to get through}

\label{lesson:HI-C416-guzaarna}



\begin{quote}
Before the new one: say the Hindi for to put to sleep, then the Hindi for a roost.
\end{quote}

\subsection*{You'll want to know: \hi{गुज़ारना}}
\textbf{\hi{गुज़ारना}} — \emph{guzārnā} — \textquotedblleft{}to spend (time); to get through\textquotedblright{}.

\begin{quote}
\hi{हमने} \hi{जंगल} \hi{में} \hi{दो} \hi{रातें} \hi{गुज़ारीं।}
\end{quote}

\begin{quote}
\emph{hamne jaṅgal meṁ do rāteṁ guzārīṁ.}
\end{quote}

\begin{quote}
We spent two nights in the forest.
\end{quote}

Past, with \hi{ने} (he or she …): \hi{उसने} \hi{वहाँ} \hi{एक} \hi{साल} \hi{गुज़ारा}, \emph{usne vahā̃ ek sāl guzārā}.

\begin{cousinweb}
\textbf{\hi{गुज़ारना}} (\emph{guzārnā}) is the causative partner of \hi{गुज़रना}, \textquotedblleft{}to pass by\textquotedblright{}, usually traced to Persian \emph{guzār}, \textquotedblleft{}passing\textquotedblright{}: time passes, and you make it pass.
\end{cousinweb}

\begin{grammarlens}[title={What you've built}]
One more word for getting around.
\end{grammarlens}

\subsection*{Your turn}
\begin{itemize}
\raggedright
  \item \emph{Say it:} \emph{guzārnā}
  \item \emph{Say it:} \emph{guzārnā}, once more
  \item \emph{Say it:} \emph{hamne jaṅgal meṁ do rāteṁ guzārīṁ.} — We spent two nights in the forest.
  \item \emph{Recall:} say the Hindi for to put to sleep, then the Hindi for a roost, then say \emph{guzārnā} again
\end{itemize}

\begin{tcolorbox}[breakable,colback=teal!4,colframe=teal!35!black,title={Before you move on}]
What does \hi{गुज़ारना} mean? (\textquotedblleft{}To spend (time); to get through\textquotedblright{}.) Say it once more.
\end{tcolorbox}
\section[\texorpdfstring{\hi{बिताना} (bitānā)}{बिताना (bitānā)}]{\hi{बिताना} (bitānā) — to pass (time)}

\label{lesson:HI-C416-bitana}



\begin{quote}
Before the new one: say the Hindi for a roost, then the Hindi for to spend (time).
\end{quote}

\subsection*{You'll want to know: \hi{बिताना}}
\textbf{\hi{बिताना}} — \emph{bitānā} — \textquotedblleft{}to pass (time)\textquotedblright{}.

\begin{quote}
\hi{हमने} \hi{छुट्टियाँ} \hi{गाँव} \hi{में} \hi{बिताईं।}
\end{quote}

\begin{quote}
\emph{hamne chuṭṭiyā̃ gā̃v meṁ bitāīṁ.}
\end{quote}

\begin{quote}
We spent the holidays in the village.
\end{quote}

Past, with \hi{ने} (he or she …): \hi{उसने} \hi{बचपन} \hi{गाँव} \hi{में} \hi{बिताया}, \emph{usne bacpan gā̃v meṁ bitāyā}.

\begin{cousinweb}
\textbf{\hi{बिताना}} (\emph{bitānā}) is the causative partner of \hi{बीतना}, \textquotedblleft{}to pass (of time)\textquotedblright{}, from Sanskrit \emph{vṛtta}, \textquotedblleft{}turned; passed\textquotedblright{} (CDIAL 12069). \hi{गुज़ारना} and \hi{बिताना} say the same thing; \hi{बिताना} is the plainer, warmer word.
\end{cousinweb}

\begin{grammarlens}[title={What you've built}]
One more word for getting around.
\end{grammarlens}

\subsection*{Your turn}
\begin{itemize}
\raggedright
  \item \emph{Say it:} \emph{bitānā}
  \item \emph{Say it:} \emph{bitānā}, once more
  \item \emph{Say it:} \emph{hamne chuṭṭiyā̃ gā̃v meṁ bitāīṁ.} — We spent the holidays in the village.
  \item \emph{Recall:} say the Hindi for a roost, then the Hindi for to spend (time), then say \emph{bitānā} again
\end{itemize}

\begin{tcolorbox}[breakable,colback=teal!4,colframe=teal!35!black,title={Before you move on}]
What does \hi{बिताना} mean? (\textquotedblleft{}To pass (time)\textquotedblright{}.) Say it once more.
\end{tcolorbox}
